\section{Introduction}
\label{sec:intro}

Score-based diffusion models learn a score network to reverse a noising process~\citep{ho2020denoising,song2021scoresde,karras2022elucidating}. Latent diffusion and related flow models generate in learned latent spaces~\citep{rombach2022highresolution,esser2024sd3}. Motivated by work on intrinsic dimension~\citep{stanczuk2024manifold,kamkari2024geometric}, we study the regime $\dlat>\dint$. Fitting the empirical score can lead to training-set replication~\citep{yi2023generalization}; collapse during reverse sampling~\citep{biroli2024dynamical} is distinct from the late-training memorization observed in some settings~\citep{bonnaire2025}. Memorization raises privacy and copyright concerns~\citep{carlini2019secret,somepalli2023diffusion,carlini2023extracting} and a basic question about what the model learns. We investigate how latent dimension affects its onset.

We ask whether widening the latent space, with everything else fixed, delays memorization, and when a spectral model can explain the delay. On the first part the evidence is clear. For MLP score networks of fixed hidden width trained on synthetic mixtures with $\dint=5$, the fraction of memorized samples after 5M steps falls from $30\%$ at $\dlat=5$ to near zero ($0.1\%$) at $\dlat=40$ over five seeds, and the learned score departs from the population score progressively later as the latent widens (\cref{sec:phenomenon}). The decline is not the nearest-neighbor test losing sensitivity in high dimension: restricted to the true five-dimensional signal subspace, where the signal-copy control retains its sensitivity, memorization still falls from $29.9\%$ to $1.2\%$ against a chance level of $0.4\%$ (\cref{tab:projection-main}). On CelebA and CIFAR-10 VAE latents, once the VAE is wide enough to preserve sample identity, the mean time to reach $5\%$ memorization increases by $3.6\times$ and $1.4\times$, respectively, over post-threshold width comparisons, while FID follows a separate quality tradeoff (\cref{sec:real}). On the second part the answer is partial: a solvable frozen-feature model identifies one mechanism for the delay (\cref{sec:theory}), but its transfer to learned representations is not established (\cref{sec:real,sec:discussion}).

Our contributions are:
\begin{itemize}
  \item \textbf{The phenomenon, with a dimension-controlled test.} We characterize the decline in memorization from $30\%$ to near zero in a fixed-width $\dlat$ sweep at known $\dint$, and test it with a projected variant of the nearest-neighbor ratio criterion. The variant is necessary: a sampler that copies signal coordinates and redraws null coordinates is invisible to the full-space test for $\dlat\ge10$, so that test alone cannot separate a real delay from a loss of sensitivity (\cref{sec:setup,sec:phenomenon}).
  \item \textbf{A spectral explanation in frozen features.} Motivated by the random-feature analysis of \citet{bonnaire2025}, we interpret measured spectra using signal, noise-dimension, sample and tail index groups (\cref{prop:fourblock-9}). Their separation and approximate edge scales are empirical and conditional, rather than a direct consequence of that theorem. Under gradient flow, modes learn in parallel and the measured ratio of signal to sample-block edges grows with $\dlat$ as the pre-activation variance $q$ falls and nonlinear feature mass shrinks. Interpreting this ratio as a memorization clock requires spectral-alignment assumptions. A control that holds $q$ fixed as $\dlat$ grows largely removes the endpoint growth ($1.09\times$ against $3.17\times$ over $\dlat=40\to240$), supporting the operating-point explanation (\cref{sec:theory}).
  \item \textbf{Real-data evidence and the transfer boundary.} CelebA and CIFAR-10 sweeps quantify memorization timing and image quality after identity preservation. Measurements in trained MLPs and a spatial-DiT pilot delimit the transfer of the frozen-feature explanation (\cref{sec:real}).
\end{itemize}

\paragraph{Relation to prior work.}
\citet{bonnaire2025} relate random-feature diffusion training to generalization and memorization timescales, and derive a limiting spectral description for Gaussian data with general clean-data covariance. \citet{george2025dsm} study Gaussian data supported on a subspace. We use these results as motivation for an anisotropic-mixture spectral approximation; finite-rank signal modes and noise averaging require additional arguments (\cref{app:operator-crosswalk}). The contribution is the operating-point interpretation and its empirical control, not a new general spectral theorem. Generative memorization is also studied through membership inference, reconstruction attacks and data-copying tests~\citep{hayes2017logan,hilprecht2019reconstruction,meehan2020datacopying}, and diffusion replication and extraction~\citep{somepalli2023diffusion,carlini2023extracting}. Our nearest-neighbor ratio follows \citet{bonnaire2025}.